% Part: computability
% Chapter: recursive-functions
% Section: other-recursions

\documentclass[../../../include/open-logic-section]{subfiles}

\begin{document}

\olfileid{cmp}{rec}{ore}
\olsection{صورت‌های دیگر بازگشت}

با استفاده از زوج‌سازی و کدگذاری دنباله‌ها می‌توانیم صورت‌های
نامتعارف‌تر (و سودمندتر) بازگشت اولیه را توجیه کنیم. برای نمونه، اغلب
سودمند است دو تابع را همزمان تعریف کنیم، مانند تعریف زیر:
\begin{align*}
h_0(\vec x, 0) & = f_0(\vec x) \\
h_1(\vec x, 0) & = f_1(\vec x) \\
h_0(\vec x, y+1) & = g_0(\vec x, y, h_0(\vec x, y), h_1(\vec x, y)) \\
h_1(\vec x, y+1) & = g_1(\vec x, y, h_0(\vec x, y), h_1(\vec x, y)).
\end{align*}
این نمونه‌ای از 
\emph{بازگشت همزمان} است. راه سودمند دیگری برای تعریف توابع آن است که
ارزش~$h(\vec x,y+1)$ را برحسب 
\emph{همهٔ} مقادیر $h(\vec x,0)$، \dots، $h(\vec x,y)$ به دست دهیم؛
مانند تعریف زیر:
\begin{align*}
h(\vec x, 0) & = f(\vec x) \\
h(\vec x, y+1) & = g(\vec x, y, \tuple{h(\vec x, 0), \dots, h(\vec x, y)}).
\end{align*}
طرحوارهٔ زیر این اندیشه را فشرده‌تر بیان می‌کند:
\[
h(\vec x, y) = g(\vec x, y, \tuple{h(\vec x, 0), \dots, h(\vec x, y-1)}),
\]
با این قرارداد که وقتی $y=0$، آخرین آرگومان~$g$ تنها دنبالهٔ تهی
است. در هر دو صورت‌بندی، اندیشه این است که هنگام محاسبهٔ «گام جانشین»،
تابع~$h$ می‌تواند از کل دنبالهٔ مقادیری استفاده کند که تاکنون محاسبه
شده‌اند. این روش 
\emph{بازگشت بر مقادیر} نامیده می‌شود. برای نمونه‌ای مشخص، می‌توان با
آن تعریف‌هایی از صورت زیر را توجیه کرد:
\begin{align*}
h(\vec x, y) & = \begin{cases}
  g(\vec x, y, h(\vec x, k(\vec x, y))) & \text{اگر $k(\vec x, y) < y$} \\
  f(\vec x) & \text{در غیر این صورت}.
\end{cases}
\end{align*}
به بیان دیگر، ارزش~$h$ در~$y$ را می‌توان برحسب ارزش~$h$ در
\emph{هر} مقدار پیشین که~$k$ به دست می‌دهد، محاسبه کرد.


\begin{prob}
  تابع باقیماندهٔ~$r(x,y)$ را با بازگشت بر مقادیر تعریف کنید. (اگر
  $x$ و~$y$ اعداد طبیعی و $y>0$ باشند، $r(x,y)$ عددی کوچک‌تر از~$y$
  است چنان‌که برای یک~$z$،
  $x=z\times y+r(x,y)$. برای معین‌بودن تعریف، قرار می‌گذاریم اگر
  $y=0$، آنگاه $r(x,0)=0$.)
\end{prob}

باید بیندیشید که چگونه می‌توان این توابع را با بازگشت اولیهٔ معمول به
دست آورد. صورت نهایی دیگری از بازگشت اولیه انعطاف‌پذیرتر است، زیرا
اجازه می‌دهد 
\emph{پارامترها} (مقادیر جانبی) را در طول راه تغییر دهیم:
\begin{align*}
h(\vec x, 0) & = f(\vec x) \\
h(\vec x, y+1) & = g(\vec x, y, h(k(\vec x), y)).
\end{align*}
این صورت را نیز می‌توان با بازگشت اولیهٔ معمول شبیه‌سازی کرد. (انجام
این کار دشوار است. برای راهنمایی، بکوشید محاسبه را با دست و از پایان
به آغاز باز کنید.)

\end{document}
